% Shared preamble for the repository diagrams. Each figure is compiled twice by build.sh: once as is
% (light theme) and once with \darktheme defined, which swaps the palette for a dark background.
\documentclass[tikz,border=4pt]{standalone}
\usepackage[T1]{fontenc}
\usepackage{lmodern}
\renewcommand{\familydefault}{\sfdefault}
\usepackage{amsmath,amssymb}
\usepackage{sansmath}
\sansmath
% Narrow boxes read better with whole words than with hyphenated ones.
\hyphenpenalty=10000
\exhyphenpenalty=10000
\usetikzlibrary{arrows.meta,positioning,calc,shapes.misc}

% Fill colors encode the role of a step: amber = test or rule, teal = forward simulation,
% purple = optimization or candidate generation, green = state or returned action, red = adverse
% outcome, blue = other processing.
\ifdefined\darktheme
  \definecolor{ink}{HTML}{E6EDF3}
  \definecolor{edge}{HTML}{8B949E}
  \definecolor{blue}{HTML}{1B2E45}
  \definecolor{teal}{HTML}{153C39}
  \definecolor{green}{HTML}{1C3823}
  \definecolor{amber}{HTML}{3F3313}
  \definecolor{red}{HTML}{442022}
  \definecolor{purple}{HTML}{30264A}
  \definecolor{grey}{HTML}{262B30}
  \definecolor{yes}{HTML}{56D364}
  \definecolor{no}{HTML}{F47067}
\else
  \definecolor{ink}{HTML}{17202A}
  \definecolor{edge}{HTML}{5D6670}
  \definecolor{blue}{HTML}{DDEBF7}
  \definecolor{teal}{HTML}{DDF2EF}
  \definecolor{green}{HTML}{E4F2E5}
  \definecolor{amber}{HTML}{FFF0CC}
  \definecolor{red}{HTML}{F8DDDD}
  \definecolor{purple}{HTML}{EAE2F5}
  \definecolor{grey}{HTML}{ECEFF1}
  \definecolor{yes}{HTML}{236B3B}
  \definecolor{no}{HTML}{A33636}
\fi

\tikzset{
  every picture/.style={color=ink, font=\small},
  box/.style={draw=edge, rounded corners=3pt, align=center, inner sep=4pt,
              text width=#1, minimum height=12mm, text=ink},
  box/.default=2.7cm,
  proc/.style={box=#1, fill=blue},   proc/.default=2.7cm,
  sim/.style={box=#1, fill=teal},    sim/.default=2.7cm,
  opt/.style={box=#1, fill=purple},  opt/.default=2.7cm,
  out/.style={box=#1, fill=green},   out/.default=2.7cm,
  bad/.style={box=#1, fill=red},     bad/.default=2.7cm,
  rule/.style={box=#1, fill=amber},  rule/.default=2.7cm,
  ask/.style={box=#1, fill=amber, chamfered rectangle, chamfered rectangle xsep=2pt,
              rounded corners=0pt, line width=0.6pt},
  ask/.default=2.7cm,
  arr/.style={-{Latex[length=2mm]}, semithick, draw=edge},
  darr/.style={arr, dashed},
  lab/.style={font=\scriptsize\bfseries, inner sep=1.5pt},
  ylab/.style={lab, text=yes}, nlab/.style={lab, text=no},
  note/.style={font=\scriptsize, text=ink, align=center},
}

\begin{document}
% How the dispatch program predicts the flow on a monitored line: the observed AC flow is the base
% point and the PTDF matrix predicts only the change caused by the controls. Illustrative values;
% the cheapest controls are used first, and curtailment only when they cannot reach the target.
\begin{tikzpicture}
\def\P#1{(#1-80)*0.32}
\def\rowh{0.55}
\foreach \v in {80,85,90,95,100,105,110}{\draw[draw=edge] ({\P{\v}},-0.15)--++(0,-0.1) node[below, font=\scriptsize, text=ink]{\v\%};}
\draw[draw=edge] (0,-0.15)--({\P{110}},-0.15);
\draw[draw=okc, dashed, thick] ({\P{95}},-0.15)--({\P{95}},4.9) node[above, font=\scriptsize, text=okc]{target 95\%};
\draw[draw=overc, dashed, thick] ({\P{100}},-0.15)--({\P{100}},4.9) node[above, font=\scriptsize, text=overc]{rating};
% waterfall: rows from top to bottom
\foreach \y/\a/\b/\c/\t in {
  4.1/80/108/blue/{observed AC flow $f^{\mathrm{AC}}_\ell$},
  3.2/99/108/purple/{redispatch, through the PTDF},
  2.3/95/99/purple/{storage, through the PTDF},
  0.5/80/95/green/{predicted flow $\hat f_\ell$}}{
  \fill[\c, draw=edge] ({\P{\a}},\y) rectangle ({\P{\b}},\y+\rowh);
  \node[anchor=east, font=\scriptsize, text=ink] at (-0.2,\y+0.5*\rowh) {\t};}
\fill[red, draw=edge] ({\P{100}},4.1) rectangle ({\P{108}},4.1+\rowh);
\node[anchor=east, font=\scriptsize, text=ink] at (-0.2,1.4+0.5*\rowh) {curtailment};
\node[anchor=west, font=\scriptsize\itshape, text=ink] at ({\P{80}+0.1},1.4+0.5*\rowh) {not needed here: it costs more than redispatch and storage};
\node[font=\scriptsize, text=ink] at ({\P{103.5}},3.2+0.5*\rowh) {$-9$};
\node[font=\scriptsize, text=ink] at ({\P{97}},2.3+0.5*\rowh) {$-4$};
\node[note, font=\small, anchor=north] at ({\P{95}},-0.75) {flow on one monitored line, in \% of its rating (illustrative values)};
% the program and its check
\node[opt=5.4cm, anchor=north west] (qp) at (11.6,4.95) {$\hat f=f^{\mathrm{AC}}+\mathrm{PTDF}\,\Delta i$\\[2pt]\scriptsize the AC solution is the base point; the linear model predicts only the change $\Delta i$ that redispatch, curtailment and storage cause};
\node[rule=5.4cm, anchor=north] (cons) at ([yshift=-4mm]qp.south) {\scriptsize each generator moves within its own ramp margin; overloads cost most, then curtailment, then relaxing a security limit; redispatch and storage are cheap};
\node[sim=5.4cm, anchor=north] (chk) at ([yshift=-4mm]cons.south) {\scriptsize the solution becomes one action, offered only if the forecast lowers the peak loading $F_H$ below that of do-nothing};
\draw[arr] (qp)--(cons); \draw[arr] (cons)--(chk);
\stagemap{3}{zoom: AC-anchored dispatch}
\end{tikzpicture}
\end{document}
